\documentclass[aspectratio=169,11pt]{beamer}
% Shared look for every deck: Metropolis with a progress bar under each frame title.
% Build with Tectonic: `make` in this directory.

\usetheme[progressbar=frametitle,numbering=fraction,sectionpage=progressbar,block=fill]{metropolis}

% Metropolis looks Fira Sans up by system name; load it from the TeX bundle instead,
% Light for body text as Metropolis intends, with lining figures for data.
\usepackage[sfdefault,light,lf]{FiraSans}
\usepackage{booktabs,tabularx}
\usepackage{tikz}
\usetikzlibrary{arrows.meta,positioning,shapes.geometric,calc}

% Metropolis's palette, with the decks' own colour names mapped onto it.
\definecolor{mDarkTeal}{HTML}{23373B}
\definecolor{mDarkBrown}{HTML}{604C38}
\definecolor{mLightBrown}{HTML}{EB811B}
\colorlet{crust}{mDarkBrown}
\definecolor{crumb}{HTML}{FAEEE2}
\colorlet{rye}{mDarkTeal}
\definecolor{flag}{HTML}{B03A2E}
\definecolor{own}{HTML}{2E7D4F}
\definecolor{muted}{HTML}{7A7F80}

% Top-align frames, so the \vfill in \source puts the source line at the bottom.
% A slimmer title bar and narrower margins than Metropolis's defaults leave room
% for the denser slides.
\makeatletter
\beamer@centeredfalse
\setlength{\metropolis@frametitle@padding}{1.4ex}
% A progress bar you can see: under every frame title and on section pages.
\setlength{\metropolis@progressinheadfoot@linewidth}{2pt}
\setlength{\metropolis@progressonsectionpage@linewidth}{1.5pt}
\setlength{\metropolis@titleseparator@linewidth}{1pt}
% Metropolis's title page is \paperheight tall, which overflows a top-aligned frame.
\setbeamertemplate{title page}{%
  \begin{minipage}[b][\textheight]{\textwidth}
    \vfill
    \usebeamertemplate*{title}
    \ifx\insertsubtitle\@empty\else\usebeamertemplate*{subtitle}\fi
    \usebeamertemplate*{title separator}
    \usebeamertemplate*{author}
    \usebeamertemplate*{institute}
    \vfill
  \end{minipage}}
\makeatother
\setbeamersize{text margin left=1.2em,text margin right=1.2em}

\setbeamertemplate{frame footer}{\insertshorttitle}
\setbeamertemplate{footline}{%
  \begin{beamercolorbox}[wd=\paperwidth,ht=2.6ex,dp=1.2ex,leftskip=1.2em,rightskip=1.2em]{footline}%
    \usebeamerfont{page number in head/foot}%
    \usebeamertemplate*{frame footer}\hfill\usebeamertemplate*{frame numbering}%
  \end{beamercolorbox}}
\setbeamerfont{page number in head/foot}{size=\tiny}
\setbeamercolor{footline}{fg=muted}

% The practical point of a slide, for the baker.
\newcommand{\takeaway}[1]{%
  \par\medskip
  \colorbox{crumb}{\parbox{\dimexpr\linewidth-2\fboxsep}{\small\textbf{\color{crust}In the kitchen:} #1}}}

% Small source line, placed at the bottom of a frame.
\newcommand{\source}[1]{%
  \vfill{\color{muted}\tiny Source: #1}}

% Status flags, mirroring the notes' conventions. Filled tags, so they read on the
% dark title bar as well as on the slide body.
\newcommand{\flagtag}[2]{\mbox{\setlength{\fboxsep}{1.2pt}\colorbox{#1}{\color{white}\scriptsize\bfseries\,#2\,}}}
\newcommand{\hypothesis}{\flagtag{flag}{hypothesis}}
\newcommand{\contested}{\flagtag{flag}{contested}}
\newcommand{\ownresult}{\flagtag{own}{own result}}
\newcommand{\untested}{\flagtag{flag}{not yet tested}}

% Source lists: one source per bullet, author and year in bold, optional group
% headings. Use in a [allowframebreaks] frame; long lists continue on "(cont.)" slides.
\newenvironment{sourcelist}{%
  \setbeamerfont{itemize/enumerate body}{size=\scriptsize}%
  \setbeamerfont{itemize item}{size=\scriptsize}%
  \scriptsize\begin{itemize}\setlength{\itemsep}{2pt}}{\end{itemize}}
\newcommand{\src}[2]{\item\textbf{#1}\enspace #2}
% \nobreak keeps a group heading on the same slide as its first source.
\newcommand{\srcgroup}[1]{\par\smallskip{\small\bfseries\color{crust}#1}\par\nobreak\vspace{-2pt}\nobreak}
\setbeamertemplate{frametitle continuation}[from second][(cont.)]

% Flowchart nodes for TikZ diagrams.
\tikzset{
  box/.style={draw=crust,fill=crumb,rounded corners=2pt,align=center,
    font=\scriptsize,inner sep=4pt,text width=#1},
  box/.default=2.6cm,
  arrow/.style={-{Stealth[length=5pt]},draw=crust,thick},
}

\author{Sourdough Science}
\institute{github.com/zrrbite/sourdough-science}
\date{}


% Local styles for this deck only.
\tikzset{
  cell/.style={draw=crust,fill=crumb,rounded corners=2pt,align=center,
    font=\footnotesize,minimum height=1.45cm,text width=2.65cm},
  axislabel/.style={font=\scriptsize,color=muted,align=right},
  colhead/.style={font=\scriptsize\bfseries,color=muted,anchor=south},
}
\newcommand{\muted}[1]{{\color{muted}#1}}
\newcolumntype{L}{>{\raggedright\arraybackslash}X}

\title[Foundations]{Preferments: The Foundations}
\subtitle{What a preferment is, and the two axes that sort them:\\
  microbiology and hydration}

\begin{document}

\begin{frame}[plain]\titlepage\end{frame}

% ---------------------------------------------------------------
\begin{frame}{What this deck covers}
  Practical questions, answered from Part~I of the notes:
  \medskip
  \begin{itemize}\setlength{\itemsep}{0.5em}
    \item What does a preferment actually do for my dough?
    \item Biga, poolish, stiff levain, liquid levain: what separates them?
    \item What does a wild culture do that commercial yeast cannot?
    \item How do I steer a levain toward a sharper or a milder sour?
    \item Why would a stiff levain survive a hot build better than a liquid one?
  \end{itemize}
  \source{chapters 01--04 of the notes}
\end{frame}

% ===============================================================
\section{A preferment and its two axes}

\begin{frame}{A preferment is dough fermented ahead of the final mix}
  \centering
  \begin{tikzpicture}
    \node[box=8cm,font=\footnotesize] (p)
      {\textbf{Preferment:} a portion of dough fermented before the final mix};
    \node[box=3.6cm,below=0.6cm of p] (gl)
      {\textbf{Gluten:} structure modified\\before the final dough};
    \node[box=3.6cm,left=0.5cm of gl] (fl)
      {\textbf{Flavour:} complexity from\\extended fermentation};
    \node[box=3.6cm,right=0.5cm of gl] (gas)
      {\textbf{Gas:} a head start\\on gas production};
    \draw[arrow] (p.south) -- (gl.north);
    \draw[arrow] (p.south -| fl.north) -- (fl.north);
    \draw[arrow] (p.south -| gas.north) -- (gas.north);
  \end{tikzpicture}
  \par\medskip\raggedright
  \begin{itemize}
    \item Every preferment does these three jobs
    \item They diverge sharply in microbiology and hydration, and so in how
      the dough behaves and what the bread becomes
  \end{itemize}
  \takeaway{choosing a preferment is choosing its organisms and its water;
    those two decide what it does to your dough and your bread.}
  \source{ch.~01, \emph{What Is a Preferment?}}
\end{frame}

% ---------------------------------------------------------------
\begin{frame}{Two independent choices define every preferment}
  \begin{columns}[T,onlytextwidth]
    \begin{column}{0.44\textwidth}
      \begin{itemize}\setlength{\itemsep}{0.3em}
        \item \textbf{Microbiology:} a monoculture of commercial
          \textit{S.~cerevisiae}, or a mixed consortium of wild yeast and
          lactic acid bacteria (LAB)
        \item \textbf{Hydration:} stiff (50--60\%) or liquid (100\%)
        \item Four types, each with its own flavour, fermentation kinetics
          and dough rheology
      \end{itemize}
    \end{column}
    \begin{column}{0.54\textwidth}
      \centering
      \begin{tikzpicture}
        \node[cell] (sl) at (0,1.6)
          {\textbf{Stiff levain}\\[1pt]\muted{sharper, vinegary\\acetic}};
        \node[cell] (ll) at (3.05,1.6)
          {\textbf{Liquid levain}\\[1pt]\muted{milder, rounder\\lactic}};
        \node[cell] (bi) at (0,0)
          {\textbf{Biga}\\[1pt]\muted{tighter, elastic\\minimal acid}};
        \node[cell] (po) at (3.05,0)
          {\textbf{Poolish}\\[1pt]\muted{very extensible\\enzymatic}};
        % hydration axis
        \draw[arrow] ($(bi.south west)+(0,-0.2)$) -- ($(po.south east)+(0,-0.2)$);
        \node[axislabel,anchor=north west] at ($(bi.south west)+(0,-0.25)$)
          {Stiff · 50--60\%};
        \node[axislabel,anchor=north east] at ($(po.south east)+(0,-0.25)$)
          {Liquid · 100\%};
        \node[font=\scriptsize\bfseries,anchor=north] at
          ($(bi.south east)!0.5!(po.south west)+(0,-0.25)$) {Hydration};
        % microbiology axis
        \draw[arrow] ($(bi.south west)+(-0.2,0)$) -- ($(sl.north west)+(-0.2,0)$);
        \node[axislabel,anchor=east] at ($(sl.west)+(-0.3,0)$)
          {Sourdough\\(wild)};
        \node[axislabel,anchor=east] at ($(bi.west)+(-0.3,0)$)
          {Commercial\\yeast};
        \node[font=\scriptsize\bfseries,anchor=south] at
          ($(sl.north west)+(-0.2,0.05)$) {Microbiology};
      \end{tikzpicture}
    \end{column}
  \end{columns}
  \takeaway{make the two decisions separately: which organisms, and how much
    water.}
  \source{ch.~02, \emph{The Four Quadrants}}
\end{frame}

% ---------------------------------------------------------------
\begin{frame}{Each quadrant has its own flavour and dough}
  \footnotesize
  \renewcommand{\arraystretch}{1.0}
  \begin{tabularx}{\textwidth}{@{}>{\bfseries\raggedright}p{1.75cm}LL@{}}
    \toprule
     & \textbf{Stiff (50--60\%)} & \textbf{Liquid (100\%)} \\
    \midrule
    Commercial yeast &
      \textbf{Biga} · Italian tradition\newline
      Long, cool ferment (12--24\,h)\newline
      Yeast aromatics: esters, aldehydes, alcohols\newline
      Minimal acidification\newline
      Tighter, more elastic gluten &
      \textbf{Poolish} · French/Polish origin\newline
      8--16\,h at room temperature\newline
      More amylase and protease (more water)\newline
      Sweeter, wheaty flavour\newline
      Very extensible dough \\
    \midrule
    Sourdough (wild) &
      \textbf{Stiff levain} · heterofermentative LAB\newline
      Lower water activity: toward acetic acid\newline
      Sharper, more vinegary tang\newline
      Wider timing window for the baker &
      \textbf{Liquid levain} · homofermentative LAB\newline
      Higher water activity: promotes lactate\newline
      Milder, rounder acidity\newline
      Handles like a poolish; full sourdough microbiology \\
    \bottomrule
  \end{tabularx}
  \takeaway{extensible dough: poolish. Elastic gluten: biga. Wider timing
    window: stiff levain.}
  \source{ch.~02, \emph{The Four Quadrants}}
\end{frame}

% ===============================================================
\section{One yeast or a consortium}

\begin{frame}{Biga and poolish run on one yeast: aroma, not acid}
  \begin{itemize}\setlength{\itemsep}{0.35em}
    \item A pure culture of \textit{S.~cerevisiae}: domesticated, single strain
    \item The long ferment is not there to grow a microbiome; it gives the
      yeast time to build flavour precursors slowly
    \item Alcohols, esters and aldehydes accumulate: nuanced aroma without
      significant acidification
    \item Some minor organic acid from the yeast's own metabolism, nothing
      near the lactic and acetic acid that LAB make
    \item The dough stays near neutral (pH 5.5--6.0), so acid does not modify
      the gluten
  \end{itemize}
  \takeaway{a long biga or poolish buys aroma, not sourness. For a sweet,
    wheaty character from enzymatic starch breakdown, the wetter poolish is
    the stronger of the two.}
  \source{ch.~03, \emph{Microbiology}; ch.~04, \emph{Acidity}}
\end{frame}

% ---------------------------------------------------------------
\begin{frame}{A levain is a stable partnership of wild yeast and bacteria}
  \begin{itemize}\setlength{\itemsep}{0.3em}
    \item \textbf{Wild yeast:} mostly wild \textit{S.\ cerevisiae} at home,
      often \textit{Kazachstania humilis} in bakeries
    \item \textbf{Lactic acid bacteria (LAB):} mostly lactobacilli, such as
      \textit{L.\ plantarum}, \textit{L.\ brevis}, \textit{F.\ sanfranciscensis}
    \item A stable association; how symbiotic is debated (Starters deck)
  \end{itemize}
  \bigskip
  \centering
  \begin{tikzpicture}[node distance=0.45cm]
    \node[box=2.6cm,minimum height=1cm] (a) {LAB make lactic\\and acetic acid};
    \node[box=2.6cm,minimum height=1cm,right=of a] (b)
      {pH drops into\\the 3.5--4.5 range};
    \node[box=2.6cm,minimum height=1cm,right=of b] (c)
      {Wild yeast, adapted\\to the acid, thrive};
    \node[box=3.8cm,minimum height=1cm,right=of c] (d)
      {Competitors suppressed:\\a stable, self-sustaining culture};
    \draw[arrow] (a) -- (b);
    \draw[arrow] (b) -- (c);
    \draw[arrow] (c) -- (d);
  \end{tikzpicture}
  \raggedright
  \takeaway{the sourness is not a side effect. The low pH is what keeps
    competing organisms out and the culture stable.}
  \source{ch.~03, \emph{Microbiology}; Landis et al.\ (2021); Carbonetto et al.\ (2020)}
\end{frame}

% ---------------------------------------------------------------
\begin{frame}{The mixed culture changes the dough in ways yeast alone cannot}
  \begin{itemize}\setlength{\itemsep}{0.15em}
    \item \textbf{Gluten is cut:} proteolytic enzymes from the LAB partially
      degrade the gluten network during fermentation
    \item \textbf{Gas is uneven:} organisms with different growth curves and
      temperature optima compete and cooperate, so gas production is less
      uniform
    \item \textbf{Gluten charge shifts:} the acid load directly modifies
      gluten charge state and hydration behaviour
    \item \contested\ Bottled acid is not the same: several studies report that
      adding acid directly \emph{firms} dough and cuts extensibility, while
      long fermentation (acid plus time-dependent proteolysis) softens it.
      The notes describe the fermentation case.
  \end{itemize}
  \takeaway{expect a levain dough to keep changing as it ferments, in ways
    biga and poolish cannot reproduce.}
  \source{ch.~03, \emph{Microbiology}; charge-state mechanism: Clarke et al.\
    (2004); contested note: ch.~25}
\end{frame}

% ===============================================================
\section{Acid and hydration}

\begin{frame}{Hydration steers which acid a levain makes}
  \centering
  \begin{tikzpicture}[
      col/.style={box=2.9cm,minimum height=1.25cm},
      wide/.style={box=3.5cm,minimum height=1.25cm},
    ]
    \def\ry{1.45}
    \def\xa{0}\def\xb{3.63}\def\xc{7.26}\def\xd{11.19}
    \foreach \x/\t in {\xa/Preferment hydration,\xb/LAB metabolism,\xc/Acid,
        \xd/Perceived · suits}
      \node[colhead] at (\x,\ry+0.65) {\t};
    % stiff levain row
    \node[col] (B) at (\xa,\ry) {\textbf{Stiff levain} · 50--60\%\\lower water activity};
    \node[col] (E) at (\xb,\ry) {Heterofermentative LAB\\shift via acetyl-CoA};
    \node[col] (H) at (\xc,\ry) {\textbf{Acetic acid}\\sharp, vinegary,\\more volatile};
    \node[wide] (J) at (\xd,\ry) {Perceived on the nose\\→ rustic, whole-grain loaves};
    % liquid levain row
    \node[col] (C) at (\xa,0) {\textbf{Liquid levain} · 100\%\\higher water activity};
    \node[col] (F) at (\xb,0) {Homofermentative LAB};
    \node[col] (I) at (\xc,0) {\textbf{Lactic acid}\\milder, rounder,\\less volatile};
    \node[wide] (K) at (\xd,0) {Perceived on the palate\\→ enriched and lighter\\breads};
    % commercial yeast row
    \node[col] (D) at (\xa,-\ry) {\textbf{Biga or poolish}\\commercial yeast};
    \node[col] (G) at (\xc,-\ry) {No meaningful\\acidification\\pH stays 5.5--6.0};
    \node[wide] (L) at (\xd,-\ry) {Gluten unmodified by acid\\→ fruity esters, sweet wheat};
    \draw[arrow] (B) -- (E); \draw[arrow] (E) -- (H); \draw[arrow] (H) -- (J);
    \draw[arrow] (C) -- (F); \draw[arrow] (F) -- (I); \draw[arrow] (I) -- (K);
    \draw[arrow] (D) -- (G); \draw[arrow] (G) -- (L);
  \end{tikzpicture}
  \raggedright
  \takeaway{hydration is one of the most important levers you have on the
    type and intensity of a levain's acid.}
  \source{ch.~04, \emph{Acidity: Lactic vs.\ Acetic and the Hydration Effect}}
\end{frame}

% ---------------------------------------------------------------
\begin{frame}{Acetic reads on the nose, lactic on the palate}
  \footnotesize
  \renewcommand{\arraystretch}{1.15}
  \begin{tabularx}{\textwidth}{@{}>{\bfseries\raggedright}p{1.9cm}LL@{}}
    \toprule
     & \textbf{Stiff levain} · 50--60\% · lower water activity
     & \textbf{Liquid levain} · 100\% · higher water activity \\
    \midrule
    LAB metabolism & Heterofermentative; shifts toward acetic acid via
      acetyl-CoA pathways when water is limiting
      & Homofermentative; produces primarily lactic acid \\
    Character & Sharper, more vinegary tang; more volatile, so perceived more
      intensely by the nose
      & Milder, rounder; less volatile, so perceived mainly on the palate \\
    Suits & Hearty, rustic loaves with whole-grain or high-extraction
      flour, where the sharpness reads as complexity
      & Enriched doughs, lighter breads, or wherever the acid should stay in
      the background and let wheat, butter or other ingredients come forward \\
    \bottomrule
  \end{tabularx}
  \takeaway{rustic whole-grain loaf: build the levain stiff. Enriched or
    lighter bread: liquid.}
  \source{ch.~04, \emph{Acidity: Lactic vs.\ Acetic and the Hydration Effect}}
\end{frame}

% ---------------------------------------------------------------
\begin{frame}{The stiff levain is also the heat-tolerant one}
  \begin{columns}[T,onlytextwidth]
    \begin{column}{0.50\textwidth}
      \begin{itemize}\setlength{\itemsep}{0.35em}
        \item Lower water activity does more than shift LAB toward acetic
          acid
        \item It makes the matrix itself protective against heat
        \item Why: irreversible protein unfolding and membrane disordering
          are both water-mediated
        \item Matters only for hot builds (ch.~11)
      \end{itemize}
    \end{column}
    \begin{column}{0.48\textwidth}
      \centering
      \begin{tikzpicture}
        \node[box=3.6cm] (s) {\textbf{Stiff levain} · 50--60\%\\lower water activity};
        \node[box=2.7cm,below left=0.7cm and -1.6cm of s] (a)
          {LAB shift toward\\acetic acid:\\\textbf{sharper}};
        \node[box=2.7cm,below right=0.7cm and -1.6cm of s] (h)
          {Matrix protects\\against heat:\\\textbf{more heat-tolerant}};
        \draw[arrow] (s) -- (a);
        \draw[arrow] (s) -- (h);
      \end{tikzpicture}
    \end{column}
  \end{columns}
  \takeaway{building levains hot? A stiff one gives the culture more
    protection, and the sharper acid comes with it.}
  \source{ch.~04, \emph{Acidity}; detail in ch.~11, \emph{Water Activity:
    The Matrix Is a Thermal Protectant}}
\end{frame}

% ===============================================================
\begin{frame}{Takeaways}
  \begin{itemize}\setlength{\itemsep}{0.3em}
    \item A preferment is dough fermented ahead: flavour, gluten change and a
      head start on gas
    \item Two independent choices define it: commercial yeast or wild culture;
      stiff (50--60\%) or liquid (100\%)
    \item Biga and poolish: one yeast, aroma without sourness; the dough stays
      at pH 5.5--6.0
    \item A levain: wild yeast plus LAB, which drop the pH to 3.5--4.5; acid
      and LAB enzymes reshape the gluten
    \item Stiff levain → acetic: sharp, on the nose, for rustic whole-grain
      loaves
    \item Liquid levain → lactic: mild, on the palate, for enriched and
      lighter breads
    \item The stiff levain is also the more heat-tolerant one, if you build
      hot
  \end{itemize}
\end{frame}

% ---------------------------------------------------------------
\begin{frame}[allowframebreaks]{Sources}
  {\scriptsize This deck draws on chapters 01--04 of the notes, which cite no literature inline. The references below are from ch.~25 and the starter note, and bear on the same ground.\par}
  \begin{sourcelist}
    \src{Clarke CI et al.\ (2004)}{Wheat sourdough fermentation: effects of time and acidification on fundamental rheological properties. \emph{Cereal Chemistry} 81(3):409--417. \muted{Source for the charge-state mechanism.}}
    \src{Thiele C, Grassl S, Gänzle MG (2004)}{Gluten hydrolysis and depolymerization during sourdough fermentation. \emph{J.\ Agric.\ Food Chem.} 52(5):1307--1314.}
    \src{Wehrle K, Grau H, Arendt EK (1997)}{Effects of lactic acid, acetic acid, and table salt on fundamental rheological properties of wheat dough. \emph{Cereal Chemistry} 74(6):739--744.}
    \src{De Vuyst L, Van Kerrebroeck S, Leroy F (2017)}{Microbial ecology and process technology of sourdough fermentation. \emph{Advances in Applied Microbiology} 100:49--160.}
    \src{Carbonetto B et al.\ (2020)}{Interactions between \textit{Kazachstania humilis} yeast species and lactic acid bacteria in sourdough. \emph{Microorganisms} 8(2):240. \muted{No mutual benefit in population terms; the pair still fermented better.}}
    \src{Landis EA et al.\ (2021)}{The diversity and function of sourdough starter microbiomes. \emph{eLife} 10:e61644. \muted{Who lives in 500 home starters.}}
  \end{sourcelist}
  \medskip
  {\scriptsize\muted{Contested evidence, ch.~25: chemical and biological acidification do not produce equivalent rheology.}\par}
\end{frame}

\end{document}
